% ============================================================
%  Expressões Numéricas — Apostila Premium | PratiqueJá
%  Engine: XeLaTeX
% ============================================================
\documentclass[12pt]{article}
\usepackage{../../pratiqueja}

% \hlm — destaque de resultado em modo-math (cor sem bold)
\newcommand{\hlm}[1]{{\color{darkblue}#1}}

\setsubject{Expressões Numéricas}

\begin{document}
\pagenumbering{arabic}
\setstretch{1.2}
\color{bodytext}

\heroheader{Expressões Numéricas}%
           {Hierarquia das Operações e Agrupadores}%
           {Matemática · Básico}

\setlength{\columnsep}{18pt}
\setlength{\columnseprule}{0.5pt}
\def\columnseprulecolor{\color{exborder}}

\begin{multicols}{2}

%─────────────────────────────────────────────────────────────
\TW{Def}{badgepurple}{Hierarquia das Operações}{A ordem em que as operações são feitas}

\regra{Ao calcular uma expressão numérica, as operações \textbf{não}
são feitas da esquerda para a direita — elas obedecem a uma
\textbf{ordem de prioridade}.}

\formula{%
\begin{tabular}{cl}
\textbf{1\textsuperscript{o}} & Potenciação e Radiciação\\[2pt]
\textbf{2\textsuperscript{o}} & Multiplicação e Divisão\\[2pt]
\textbf{3\textsuperscript{o}} & Adição e Subtração
\end{tabular}}

\nota{Operações de \textbf{mesma prioridade} são resolvidas da
\textbf{esquerda para a direita}.}

%─────────────────────────────────────────────────────────────
\TW{$(\,)$}{darkblue}{Parênteses, Colchetes e Chaves}{Agrupadores alteram a prioridade}

\regra{Com agrupadores, resolva \textbf{de dentro para fora}:
primeiro os \textbf{parênteses} $(\,)$, depois os \textbf{colchetes}
$[\,]$ e por fim as \textbf{chaves} $\{\,\}$. O que está dentro tem
prioridade máxima.}

\SH{Exemplo}
\exemp{%
  $\{3 \times [2 + (4 - 1)]\} + 5$\\
  $= \{3 \times [2 + 3]\} + 5$\\
  $= \{3 \times 5\} + 5$\\
  $= 15 + 5 = \hlm{20}$%
}

%─────────────────────────────────────────────────────────────
\TW{$\pm\times\div$}{darkblue}{Expressões com as 4 Operações}{Aplicando a hierarquia correta}

\SH{Exemplo 1 — sem agrupadores}
\exemp{%
  $8 + 6 \div 2 - 3 \times 2$\\
  $= 8 + 3 - 6$ \quad{\footnotesize\itshape(÷ e × primeiro)}\\
  $= 11 - 6 = \hlm{5}$%
}

\SH{Exemplo 2 — com parênteses}
\exemp{%
  $(8 + 6) \div 2 - 3 \times 2$\\
  $= 14 \div 2 - 3 \times 2$\\
  $= 7 - 6 = \hlm{1}$%
}

\obs{Compare: o mesmo conjunto de números dá resultados diferentes
($5$ e $1$) só pela presença dos parênteses.}

%─────────────────────────────────────────────────────────────
\TW{$a^n$}{darkblue}{Expressões com Potenciação}{Potências têm prioridade máxima}

\SH{Exemplo 1}
\exemp{%
  $2^3 + 4 \times 2 - 5$\\
  $= 8 + 4 \times 2 - 5$ \quad{\footnotesize\itshape(potência primeiro)}\\
  $= 8 + 8 - 5 = \hlm{11}$%
}

\SH{Exemplo 2 — agrupadores e potências}
\exemp{%
  $[3^2 - (4 + 1)] \times 2$\\
  $= [9 - 5] \times 2$\\
  $= 4 \times 2 = \hlm{8}$%
}

%─────────────────────────────────────────────────────────────
\TW{Ex}{badgegreen}{Exemplos Resolvidos}{Exercícios completos com resolução}

\SH{Exercício 1}
\exemp{%
  Calcule: $\{[(12 \div 4) + 3^2] \times 2\} - 10$\\[2pt]
  $= \{[3 + 9] \times 2\} - 10$ \quad{\footnotesize\itshape(parênt. e potência)}\\
  $= \{12 \times 2\} - 10$\\
  $= 24 - 10 = \hlm{14}$%
}

\SH{Exercício 2}
\exemp{%
  Calcule: $5 + 3 \times (2^2 - 1) \div 3$\\[2pt]
  $= 5 + 3 \times (4 - 1) \div 3$\\
  $= 5 + 3 \times 3 \div 3$\\
  $= 5 + 9 \div 3$\\
  $= 5 + 3 = \hlm{8}$%
}

\nota{\textbf{Dica de revisão:} ao terminar, refaça os passos
conferindo se cada etapa respeitou a prioridade e os agrupadores
de dentro para fora.}

\end{multicols}

\end{document}
